\documentclass[11pt,a4paper]{article}

\usepackage[a4paper,margin=1in]{geometry}
\usepackage[T1]{fontenc}
\usepackage{newtxtext,newtxmath}
\usepackage{microtype}
\usepackage{xurl}
\usepackage{hyperref}
\usepackage{tabularx}
\usepackage{enumitem}
\usepackage{amsmath}
\usepackage{amssymb}

\hypersetup{
  colorlinks=true,
  linkcolor=black,
  urlcolor=blue,
  citecolor=black
}

\setlist[itemize]{topsep=4pt,itemsep=2pt}
\setlist[enumerate]{topsep=4pt,itemsep=2pt}

\title{\textbf{SQI Semiprime Seed-Wave Breakthrough Lock}}
\author{Kevin Robinson \\ Tessaris AI}
\date{May 2026}

\begin{document}

\maketitle

\begin{abstract}
This lock records the first observed learning-feedback improvement in the SQI semiprime
factor search harness. The system learned a successful Rho/Brent seed and constant pair
from previous collapse traces, promoted it to a first-priority symbolic branch, and reduced
hard spread semiprime factorisation times from hundreds/thousands of milliseconds to
single-digit milliseconds on the repeated hard workload.
\end{abstract}

\section{Lock Status}

\begin{tabularx}{\textwidth}{@{}lX@{}}
\textbf{Lock ID} & SQI-SEED-WAVE-LEARNING-LOCK-0003 \\
\textbf{Status} & LOCKED --- first learned seed-wave feedback improvement observed \\
\textbf{Maintainer} & Tessaris AI \\
\textbf{Author} & Kevin Robinson \\
\end{tabularx}

\section{Purpose}

This document records the first meaningful SQI learning-feedback result in the prime factor
search programme.

The goal was not merely to factor a number once. The goal was to test whether SQI could:

\begin{enumerate}
  \item run semiprime stress tests;
  \item detect successful factor-search branches;
  \item extract seed and constant priors from successful collapse traces;
  \item feed those priors back into future runs;
  \item improve future factor-search performance.
\end{enumerate}

\section{Implemented Components}

The current system includes:

\begin{itemize}
  \item \url{backend/modules/sqi/factor/sqi_prime_factorizer.py}
  \item \url{backend/modules/sqi/factor/semiprime_lab.py}
  \item \url{backend/modules/sqi/factor/seed_wave_learner.py}
  \item \url{backend/modules/sqi/factor/seed_wave_profiler.py}
  \item \url{backend/modules/sqi/factor/sqi_factor_cli.py}
\end{itemize}

Associated tests:

\begin{itemize}
  \item \url{backend/tests/test_sqi_prime_factorizer.py}
  \item \url{backend/tests/test_sqi_prime_factor_benchmark_runner.py}
  \item \url{backend/tests/test_sqi_factor_cli.py}
  \item \url{backend/tests/test_sqi_semiprime_lab.py}
  \item \url{backend/tests/test_sqi_seed_wave_learner.py}
  \item \url{backend/tests/test_sqi_seed_wave_profiler.py}
\end{itemize}

\section{Validation}

Current test command:

\begin{verbatim}
PYTHONPATH=. python -m pytest -q \
  backend/tests/test_sqi_prime_factorizer.py \
  backend/tests/test_sqi_prime_factor_benchmark_runner.py \
  backend/tests/test_sqi_factor_cli.py \
  backend/tests/test_sqi_semiprime_lab.py \
  backend/tests/test_sqi_seed_wave_learner.py \
  backend/tests/test_sqi_seed_wave_profiler.py
\end{verbatim}

Observed result:

\begin{verbatim}
28 passed, 1 warning
\end{verbatim}

The warning is an existing repository-wide container import deprecation warning. It is not
a failure of the SQI factor-search harness.

\section{Experimental Setup}

The hard semiprime lab was run with:

\begin{verbatim}
PYTHONPATH=. python backend/modules/sqi/factor/semiprime_lab.py \
  --per-tier 1 \
  --close-bits 44,48 \
  --spread-pairs 32x56,36x60 \
  --json
\end{verbatim}

This generated four cases:

\begin{itemize}
  \item close \(44 \times 44\)-bit semiprime;
  \item close \(48 \times 48\)-bit semiprime;
  \item spread \(32 \times 56\)-bit semiprime;
  \item spread \(36 \times 60\)-bit semiprime.
\end{itemize}

\section{Learned Seed Prior}

The seed-wave learner extracted the following useful prior from the hard spread bucket:

\begin{verbatim}
seed = 15
c    = 14
branch = 13
\end{verbatim}

This was stored under the \texttt{spread\_64\_95} bucket in:

\begin{verbatim}
benchmarks/sqi_semiprime_lab/seed_wave_priors.json
\end{verbatim}

\section{Profiler Evidence}

The profiler showed that the hard \(36 \times 60\)-bit spread semiprime could be factored
quickly using the learned Brent seed branch:

\begin{tabularx}{\textwidth}{@{}lX@{}}
\textbf{Method} & Brent Rho seed-wave branch \\
\textbf{Seed} & 15 \\
\textbf{Constant \(c\)} & 14 \\
\textbf{Checks} & 25,726 \\
\textbf{Observed time} & approximately 5.8 ms \\
\textbf{Factor} & 57,925,206,503 \\
\end{tabularx}

The same seed/constant pair was also the fastest observed Pollard branch by checks, though
Brent was faster in wall-clock time for this case.

\section{Before and After}

Before learned-first feedback:

\begin{tabularx}{\textwidth}{@{}l l r l@{}}
\textbf{Case} & \textbf{Method} & \textbf{Approx. time} & \textbf{Verified} \\
\hline
spread \(32 \times 56\) & Pollard Rho & 216 ms & yes \\
spread \(36 \times 60\) & Brent seed-wave fallback & 1112 ms & yes \\
\end{tabularx}

After learned-first seed-wave feedback:

\begin{tabularx}{\textwidth}{@{}l l r l@{}}
\textbf{Case} & \textbf{Method} & \textbf{Approx. time} & \textbf{Verified} \\
\hline
spread \(32 \times 56\) & learned-first Brent seed-wave & 9.0 ms & yes \\
spread \(36 \times 60\) & learned-first Brent seed-wave & 6.0 ms & yes \\
\end{tabularx}

\section{Observed Hard-Run Summary}

The final observed hard run produced:

\begin{verbatim}
route_correct_family_count: 4
route_correct_family_rate: 1.0
sqi_solved: 4
sqi_solved_rate: 1.0
total_cases: 4
winner_counts:
  fermat_only: 2
  sqi: 2
\end{verbatim}

The final case-level result was:

\begin{tabularx}{\textwidth}{@{}l l r l l@{}}
\textbf{Case} & \textbf{Method} & \textbf{SQI ms} & \textbf{Verified} & \textbf{Winner} \\
\hline
close \(44 \times 44\) & Fermat & 0.055 & yes & Fermat-only \\
close \(48 \times 48\) & Fermat & 0.049 & yes & Fermat-only \\
spread \(32 \times 56\) & learned-first Brent seed-wave & 9.055 & yes & SQI \\
spread \(36 \times 60\) & learned-first Brent seed-wave & 6.002 & yes & SQI \\
\end{tabularx}

\section{Scientific Claim}

The safe scientific claim is:

\begin{quote}
SQI has demonstrated a closed symbolic factor-search learning loop: it can run semiprime
stress tests, identify successful Rho/Brent seed branches, store those branches as priors,
and reuse them as first-priority symbolic branches to materially improve repeated hard
spread-semiprime factorisation performance.
\end{quote}

This is stronger than a fixed factorisation harness because the system now adapts using
collapse history.

\section{Important Caveat}

This does not prove that SQI has cracked prime numbers in general.

It does not prove:

\begin{itemize}
  \item RSA-breaking capability;
  \item physical quantum speedup;
  \item replacement of Shor's algorithm;
  \item general prediction of prime factors without computation;
  \item superiority on unseen arbitrary semiprimes.
\end{itemize}

The result is currently a learned-feedback improvement on a repeated hard workload family.

\section{Next Scientific Test}

The next required test is generalisation:

\begin{quote}
Can learned seed-wave priors improve performance on unseen semiprimes from the same
bit-size bucket, rather than only on a repeated case?
\end{quote}

The next build should therefore add:

\begin{itemize}
  \item train/test split for semiprime labs;
  \item prior training on one batch;
  \item evaluation on unseen semiprimes;
  \item statistical comparison against Brent-only and Pollard-only;
  \item confidence table over multiple random seeds.
\end{itemize}

\section{Next Engineering Target}

Create:

\begin{verbatim}
backend/modules/sqi/factor/semiprime_generalization_lab.py
\end{verbatim}

The lab should run:

\begin{enumerate}
  \item training semiprime batch;
  \item profiler/learner extraction;
  \item unseen evaluation batch;
  \item baseline comparison;
  \item final report.
\end{enumerate}

\bigskip

\noindent\textbf{Lock ID:} SQI-SEED-WAVE-LEARNING-LOCK-0003

\noindent\textbf{Status:} LOCKED --- first learned seed-wave feedback improvement observed

\noindent\textbf{Maintainer:} Tessaris AI

\noindent\textbf{Author:} Kevin Robinson

\end{document}
